
\section*{Part II. Applications of optics \hfill 9 / 15}
\subsection*{1. Parameter estimation and Fisher information}
A probe accumulates an unknown phase \(\phi\) through
\(\ket{\psi_\phi}=e^{-i\phi G}\ket{\psi_0}\).
For a chosen measurement with outcomes \(x\),
classical Fisher information is
\[
F(\phi)=\sum_xp(x|\phi)
\left[\partial_\phi\ln p(x|\phi)\right]^2.
\]
For \(M\) independent repetitions, an unbiased estimator
obeys the Cram\'er--Rao bound
\(\operatorname{Var}(\tilde\phi)\ge1/[M F(\phi)]\).
Optimising over physically allowed measurements gives
the quantum Fisher information \(F_Q\):
\[
\boxed{\operatorname{Var}(\tilde\phi)
\ge\frac{1}{M F_Q},\qquad
F_Q=4(\Delta G)^2\quad\text{for a pure probe}.}
\]
If phase is generated by photon number \(G=\hat n\), coherent
light has \(F_Q=4\bar n\). The standard deviation therefore
scales as \(1/\sqrt{\bar n}\) (shot-noise scaling, up to
constant factors). Carefully chosen nonclassical states can
offer better scaling under specified resource and loss assumptions.

\subsection*{2. Interferometers convert phase to intensity}
A Mach--Zehnder interferometer uses two beam splitters separated
by a relative phase. In a balanced convention,
a photon entering one input has output probabilities
\[
P_1=\cos^2(\Delta\phi/2),\qquad
P_2=\sin^2(\Delta\phi/2)
\]
up to which port is named 1 or 2. The phase shift is invisible
to direct intensity measurement on a single path until the
two arms recombine. A dark output at \(\Delta\phi=0\)
becomes bright when the relative phase changes.

\subsection*{3. Why squeezing helps}
A coherent state has equal vacuum noise in every
quadrature. If the measured quadrature is squeezed below
vacuum noise, the same signal displacement can be resolved
with smaller fluctuations. This underlies squeezed-light
enhancement of precision interferometers, including
gravitational-wave detectors. Real gain depends on
optical loss, mode matching and measurement angle.

\begin{methodbox}
\textbf{Typical derivation:} identify the encoded generator \(G\);
calculate \(\langle G\rangle,\langle G^2\rangle\);
find \(F_Q\), then the bound for \(M\) repetitions.
For a concrete detector, compute \(p(x|\phi)\) first and compare
its classical \(F\) to \(F_Q\).
\end{methodbox}
\src{W6L3 pp.~5--37; W3L2 on squeezing; W2L2 on beam splitters.}
